% Appendix. The main text is self-contained; this is supporting detail.

\chapter{Standard background derivations}
\label{app:background}

This appendix records the established results used in Section~\ref{sec:m-prelim}.
They are not contributions of this thesis; they are collected here so the main text
can state them briefly.

The trajectory information bound (Eq.~\eqref{eq:traj-mi}) is derived in the main text
(Section~\ref{sec:m-policy}); we record only the contrastive bounding properties and
the per-step chain rule here.

\section{Bounding properties of the contrastive estimators}
\label{app:bounds}

Recall $\mathrm{sPCE}_L$ from Eq.~\eqref{eq:spce} and $\mathrm{sNMC}_L$, which uses
the same numerator but divides the denominator sum by $L$ over the contrastives
$\theta_{1:L}$ only.

\paragraph{sPCE is a lower bound, capped at $\log(L+1)$.}
The denominator of sPCE is an average over $L+1$ likelihoods that estimates the
same marginal $p(y_t\mid\xi_t,h_{t-1})$, but it includes the data-generating
$\theta_0$, whose likelihood is systematically large for the $y_t$ it produced.
This biases the denominator up and the ratio down, so
$\mathbb{E}[\mathrm{sPCE}_L]\le I_{h_{t-1}}(\xi_t)$. Equivalently, sPCE is a
classification bound: among the $L+1$ exchangeable draws, identifying which
generated $y_t$ has success probability at most one, so
$\mathbb{E}[\mathrm{sPCE}_L]=\log(L+1)-H(\text{index}\mid\theta_{0:L},y_t)\le\log(L+1)$.

\paragraph{sNMC is an upper bound.}
The denominator $\frac1L\sum_{\ell=1}^L p(y_t\mid\theta_\ell,\xi_t)$ is an unbiased
estimate of the marginal. By Jensen's inequality the expected log of an unbiased
estimate underestimates the log of the mean, so the log-marginal is underestimated
and the ratio overestimated, giving $\mathbb{E}[\mathrm{sNMC}_L]\ge I_{h_{t-1}}(\xi_t)$.

Both biases are $O(1/L)$, so $[\mathrm{sPCE}_L,\mathrm{sNMC}_L]$ brackets the true
EIG and tightens as $L\to\infty$.

\section{Per-step chain rule}
\label{app:chain}

Let $S_t$ be the sPCE evaluated on the length-$t$ prefix and $S_0=0$. The per-step
increment $\Delta S_t=S_t-S_{t-1}$ telescopes, $\sum_{t=1}^{T}\Delta S_t=S_T$,
recovering the trajectory bound. Each $\Delta S_t$ weights the contrastive
denominator by the running posterior, so it measures the information acquired at
step $t$ within the operative posterior rather than from a fresh prior estimate.

\chapter{Linearised EIG and the escape threshold}
\label{app:lindley}

\section{The linearised information gain equals $\Lambda$}
The main text calls $\Lambda$ the Gaussian-linearised information gain; we justify the
name. For a Gaussian prior $p(\theta)=\mathcal N(\mu_0,\Sigma_0)$ and a linear-Gaussian
observation $y=\mu_0+J(\theta-\theta_0)+\varepsilon$ with $\varepsilon\sim\mathcal N(0,R)$,
the posterior is Gaussian with covariance $\Sigma_{\mathrm{post}}=(\Sigma_0^{-1}+F)^{-1}$,
$F=J^\top R^{-1}J$, which does not depend on the realised $y$. The information gain is the
drop in differential entropy,
\begin{align}
\mathrm{EIG}(\xi)
&= H[p(\theta)]-\mathbb{E}_y\,H[p(\theta\mid y,\xi)]
= \tfrac12\log\!\big((2\pi e)^d\det\Sigma_0\big)-\tfrac12\log\!\big((2\pi e)^d\det\Sigma_{\mathrm{post}}\big)\nonumber\\
&= \tfrac12\log\frac{\det\Sigma_0}{\det\Sigma_{\mathrm{post}}}
= \tfrac12\log\det\!\big(\Sigma_0(\Sigma_0^{-1}+F)\big)
= \tfrac12\log\det(I+\Sigma_0 F)
= \tfrac12\sum_i\log(1+\lambda_i)=\Lambda. \label{eq:eig-lambda}
\end{align}
So $\Lambda$ is exactly the EIG of the linearised experiment (equivalently the Bayesian
D-optimal value). For nonlinear models the same expression is the leading term of the
entropy-based EIG under a Gauss--Newton (Laplace) approximation, with $\Sigma_{\mathrm{post}}$
evaluated at the posterior mode. Combined with the escape result below, this gives the
headline relation $L_{\mathrm{escape}}\sim e^{\Lambda}=e^{\mathrm{EIG}}$: the contrastive
budget needed to escape saturation grows as the exponential of the experiment's
information.

\paragraph{Interpretation.} Each $\lambda_i$ is a per-axis signal-to-noise ratio. It is
an eigenvalue of the prior-scaled Fisher information $\Sigma_0^{1/2}F\Sigma_0^{1/2}$,
comparing how much the experiment informs direction $i$ (through $F$) against the prior
uncertainty there (through $\Sigma_0$). A well-identified direction has $\lambda_i\gg1$
and contributes $\tfrac12\log(1+\lambda_i)\approx\tfrac12\log\lambda_i$ nats; a
barely-informed direction has $\lambda_i\ll1$ and contributes almost nothing. $\Lambda$
is the sum of these per-axis contributions, and $e^{\Lambda}=\prod_i(1+\lambda_i)^{1/2}$
is the factor by which the posterior volume shrinks below the prior, which is exactly
the contrastive budget needed to escape saturation.

\section{The expected importance weight and the escape threshold}
We derive in full the scaling $\mathbb{E}[r_\ell]\sim e^{-\Lambda}$ stated in
Eq.~\eqref{eq:escape-scale}, via the marginal-likelihood route. Take a Gaussian prior
$p(\theta)=\mathcal N(\mu_0,\Sigma_0)$ and observation $y=\mu(\theta,\xi)+\varepsilon$
with $\varepsilon\sim\mathcal N(0,R)$. The Fisher information is $F=J^\top R^{-1}J$
with $J=\partial\mu/\partial\theta|_{\theta_0}$, and the posterior precision is
$\Sigma_{\mathrm{post}}^{-1}=\Sigma_0^{-1}+F$.

\paragraph{Step 1: Laplace approximation to the marginal likelihood.}
Expand the log-likelihood to second order about $\theta_0$,
\[
\log p(y\mid\theta,\xi) \approx \log p(y\mid\theta_0,\xi) + (\theta-\theta_0)^\top s - \tfrac12(\theta-\theta_0)^\top F(\theta-\theta_0),
\]
where $s=\nabla_\theta\log p(y\mid\theta,\xi)|_{\theta_0}$ is the score and the
curvature is the observed information $H=-\nabla_\theta^2\log p(y\mid\theta,\xi)|_{\theta_0}$;
under Gaussian noise $H=J^\top R^{-1}J=F$ exactly, and for general likelihoods
$\mathbb{E}_y[H]=F$, so we use $F$ throughout. Adding the Gaussian log-prior and
completing the square,
\[
\log\bigl[p(y\mid\theta,\xi)\,p(\theta)\bigr] \approx \log\bigl[p(y\mid\theta_0,\xi)\,p(\theta_0)\bigr] + \tfrac12 g^\top\Sigma_{\mathrm{post}}\,g - \tfrac12(\theta-\hat\theta)^\top\Sigma_{\mathrm{post}}^{-1}(\theta-\hat\theta),
\]
where $g=s-\Sigma_0^{-1}(\theta_0-\mu_0)$ is the log-posterior gradient at $\theta_0$
and $\hat\theta=\theta_0+\Sigma_{\mathrm{post}}\,g$ is the posterior mode. The
unnormalised posterior is thus approximately Gaussian with covariance
$\Sigma_{\mathrm{post}}$; integrating the Gaussian kernel
($\int = (2\pi)^{d/2}\det(\Sigma_{\mathrm{post}})^{1/2}$),
\[
p(y\mid\xi) \approx p(y\mid\theta_0,\xi)\,p(\theta_0)\,e^{\frac12 g^\top\Sigma_{\mathrm{post}}\,g}\,(2\pi)^{d/2}\det(\Sigma_{\mathrm{post}})^{1/2}.
\]

\paragraph{Step 2: Expected importance weight.}
From Eq.~\eqref{eq:Er}, $\mathbb{E}[r_\ell]=\int p(y\mid\theta,\xi)\,p(\theta)\,\ud\theta\,/\,p(y\mid\theta_0,\xi)$.
Substituting the Laplace marginal and the Gaussian prior density
$p(\theta_0)=(2\pi)^{-d/2}\det(\Sigma_0)^{-1/2}\exp(-\tfrac12(\theta_0-\mu_0)^\top\Sigma_0^{-1}(\theta_0-\mu_0))$,
\[
\mathbb{E}[r_\ell] \approx \underbrace{\frac{\det(\Sigma_{\mathrm{post}})^{1/2}}{\det(\Sigma_0)^{1/2}}}_{\text{(1) volume ratio}} \cdot \underbrace{\exp\!\Bigl(-\tfrac12(\theta_0-\mu_0)^\top\Sigma_0^{-1}(\theta_0-\mu_0)\Bigr)}_{\text{(2) prior centering}} \cdot \underbrace{e^{\frac12 g^\top\Sigma_{\mathrm{post}}\,g}}_{\text{(3) mode shift}}.
\]
Factor (1) is the exponentially small volume ratio that drives saturation. Factor (2)
involves a Mahalanobis distance that is $\chi^2_d$ with mean $d$, so it is
$e^{-\Theta(d)}$ (at the mean exponent, $e^{-d/2}$; in exact expectation
$\mathbb{E}[e^{-\frac12\chi^2_d}]=2^{-d/2}$). Factor (3) is $\ge1$ with exponent
$O(d)$. Factors (2) and (3) push in opposite directions, contribute $O(d)$ to the
exponent, and do not affect the exponential scaling in $\Lambda$. Retaining the
leading volume ratio,
\[
\mathbb{E}[r_\ell]\sim\frac{\det(\Sigma_{\mathrm{post}})^{1/2}}{\det(\Sigma_0)^{1/2}}=\frac{1}{\det(I+\Sigma_0 F)^{1/2}}=\prod_i(1+\lambda_i)^{-1/2}=e^{-\Lambda}.
\]

\paragraph{Step 3: Eigenvalue identity.}
$\Sigma_0 F$ is similar to the symmetric positive semidefinite matrix
$M:=\Sigma_0^{1/2}F\Sigma_0^{1/2}$ through $P=\Sigma_0^{1/2}$,
\[
P\,M\,P^{-1}=\Sigma_0^{1/2}\bigl(\Sigma_0^{1/2}F\Sigma_0^{1/2}\bigr)\Sigma_0^{-1/2}=\Sigma_0 F .
\]
Similar matrices share eigenvalues, and $M$ is symmetric PSD, since for any $v$,
$v^\top M v=(\Sigma_0^{1/2}v)^\top F(\Sigma_0^{1/2}v)\ge0$; its eigenvalues
$\lambda_i\ge0$ are therefore real. Hence
\[
\det(I+\Sigma_0 F)=\det(I+M)=\prod_i(1+\lambda_i),
\qquad \Lambda=\tfrac12\sum_i\log(1+\lambda_i).
\]

\paragraph{Step 4: Escape threshold.}
The expected contrastive sum is $\mathbb{E}[S]=L\,e^{-\Lambda}$. By Markov's
inequality $\Pr(S\ge1)\le\mathbb{E}[S]=L\,e^{-\Lambda}$, so saturation persists while
$\mathbb{E}[S]\ll1$, i.e.\ $L\ll e^{\Lambda}$, and escape requires
$L\gtrsim e^{\Lambda}$. $\square$

\paragraph{Correction factors.}
The combined effect of factors (2) and (3) shifts $L_{\mathrm{escape},50}$ by a
multiplicative factor $\sim e^{O(d)}$, large in absolute terms for $d=12$ (ASM1) but
subexponential relative to $e^{\Lambda}$ when $\Lambda\gg d$. Table~\ref{tab:escape-full}
reports the full twelve-system measurements; the structural mechanisms that push the
\emph{empirical} threshold beyond this Laplace prediction are derived in the main text
(Section~\ref{sec:m-mechanisms}).

\begin{table}[h]
\centering\small
\caption[Full saturation thresholds]{Empirical escape thresholds at random designs,
sorted ascending, with the Laplace reference $e^\Lambda$ where available. Stiff and
Hill-kinetic systems exceed $e^\Lambda$ by one to three orders of magnitude.}
\label{tab:escape-full}
\begin{tabular}{lccrrl}
\toprule
System & $d_\theta$ & $d_\xi$ & $\Lambda$ & $e^{\Lambda}$ & $L_{\mathrm{escape},50}$ \\
\midrule
PK 8-comp & 20 & 1 & $0.05$ & $1.1$ & ${\le}10$ \\
PK 1-comp & 4 & 1 & $1.02$ & $2.8$ & ${\sim}10$ \\
PI3K      & 6 & 2 & $1.49$ & $4.4$ & ${\le}10$ \\
Monod     & 2 & 1 & $4.30$ & $74$ & ${\sim}30$ \\
CES       & 4 & 6 & -- & -- & ${\sim}10^4$ (static) \\
Goodwin   & 5 & 1 & $4.40$ & $81$ & ${\sim}10^4$ \\
ASM1      & 12 & 2 & -- & -- & ${\sim}10^4$ (stiff) \\
SEIR R=3  & 9 & 9 & $8.64$ & $5{,}648$ & ${>}10^6$ \\
SEIR R=8  & 14 & 24 & $21.26$ & $1.7{\times}10^9$ & ${>}10^6$ \\
\bottomrule
\end{tabular}
\end{table}

\chapter{PACE: consistency, decomposition, and the Fisher limit}
\label{app:pace}

\section{Consistency}
Theorem~\ref{thm:consistency} is the standard two-stage self-normalized importance
sampling result. Support inclusion makes the weighted averages in the numerator and
denominator of Eq.~\eqref{eq:pace} converge almost surely to the corresponding
posterior expectations; finite chi-square divergence gives the $O(1/J)$ bias rate
through the SNIS central limit theorem. With a perfect proposal the weights are the
constant of Eq.~\eqref{eq:widentity} and the estimator is exact in the weights.

\section{Finite variance and the likelihood-factor cancellation}
\label{app:chi2}
The $O(1/J)$ rate in Theorem~\ref{thm:consistency} requires the chi-square divergence
$\chi^2(p(\cdot\mid h)\,\|\,q_\phi(\cdot\mid h))=\mathbb{E}_{q_\phi}\!\big[(p(\theta\mid h)/q_\phi(\theta\mid h))^2\big]-1$
to be finite, which holds for any NPE that covers the posterior. It is worth seeing why
PACE's weight has this property when the naive alternative does not. Estimating the
per-step gain by drawing contrastives from $q_\phi$ but \emph{targeting the prior
marginal}, as a prior-contrastive denominator does, would weight each draw by
$p(\theta_\ell)/q_\phi(\theta_\ell\mid h)$, whose divergence is
\begin{equation}
\label{eq:chi2-prior}
\chi^2\!\big(p(\theta)\,\big\|\,q_\phi(\cdot\mid h)\big)=\int\frac{p(\theta)^2}{q_\phi(\theta\mid h)}\,\ud\theta-1 .
\end{equation}
Since $q_\phi(\cdot\mid h)$ approximates the concentrated posterior, it has negligible
mass in the prior-dominated tails where $p(\theta)$ does not vanish, so the integrand
$p(\theta)^2/q_\phi(\theta\mid h)$ diverges and this $\chi^2$ is unbounded. PACE's weight
is instead
\[
w_j=\frac{p(\theta_j)\,p(h\mid\theta_j)}{q_\phi(\theta_j\mid h)}\;\propto\;\frac{p(\theta_j\mid h)}{q_\phi(\theta_j\mid h)} ,
\]
where the likelihood factor $p(h\mid\theta_j)$ is small exactly in those prior tails and
cancels the divergence, leaving the posterior-to-proposal ratio whose $\chi^2$ is finite
for a covering NPE. This cancellation is what makes the cost transfer of
Section~\ref{sec:m-costtransfer} work: the same likelihood factor that retargets the
estimator from the prior to the posterior is what keeps its variance finite.

\section{Decomposition}
Substituting $\tilde w_k=w_k/W$ with $W=\sum_k w_k$ into Eq.~\eqref{eq:pace} and
factoring $p(y_{t,j}\mid\theta_j,\xi_t)$ out of the inner sum using
$r_{k\leftarrow j}=p(y_{t,j}\mid\theta_k,\xi_t)/p(y_{t,j}\mid\theta_j,\xi_t)$ gives,
for the $j$-th term, $\log W-\log w_j-\log[1+\Psi_j(\xi_t)]$ with
$\Psi_j(\xi_t)=\sum_{k\ne j}(w_k/w_j)r_{k\leftarrow j}$. Taking the weighted sum and
using $\sum_j\tilde w_j\log w_j=-H(\tilde w)+\log W$ yields Eq.~\eqref{eq:decomp},
$\hat I^{\mathrm{PACE}}=H(\tilde w)-\bar\Phi(\xi_t)$. The collapsed-weight limit of
Eq.~\eqref{eq:lr-limit} follows as in Section~\ref{sec:m-decomp}.

\section{The Fisher limit}
We expand the collapsed-weight criterion of Section~\ref{sec:m-collapse}. When a single
competitor $\theta_{k^\ast}=\arg\max_{k\ne\mathrm{dom}}w_k$ dominates the remaining
non-dominant particles, the design-dependent term reduces to one summand,
\begin{equation}
\label{eq:pairwise}
\bar\Phi(\xi_t)\approx\frac{w_{k^\ast}}{w_{\mathrm{dom}}}\,r_{k^\ast\leftarrow\mathrm{dom}}(\xi_t)
=\frac{w_{k^\ast}}{w_{\mathrm{dom}}}\,\frac{p(y_{t,\mathrm{dom}}\mid\theta_{k^\ast},\xi_t)}{p(y_{t,\mathrm{dom}}\mid\theta_{\mathrm{dom}},\xi_t)} .
\end{equation}
The prefactor $w_{k^\ast}/w_{\mathrm{dom}}$ is constant in $\xi_t$, so minimising
$\bar\Phi$ is maximising the pairwise log likelihood-ratio
$\log[p(y_{t,\mathrm{dom}}\mid\theta_{\mathrm{dom}},\xi_t)/p(y_{t,\mathrm{dom}}\mid\theta_{k^\ast},\xi_t)]$.
Its expectation over $y_{t,\mathrm{dom}}\sim p(y_t\mid\theta_{\mathrm{dom}},\xi_t)$ is the
Kullback-Leibler divergence
\[
\mathbb{E}_{y_{t,\mathrm{dom}}}\!\left[\log\frac{p(y_{t,\mathrm{dom}}\mid\theta_{\mathrm{dom}},\xi_t)}{p(y_{t,\mathrm{dom}}\mid\theta_{k^\ast},\xi_t)}\right]
= D_{\mathrm{KL}}\!\big(p(y_t\mid\theta_{\mathrm{dom}},\xi_t)\,\|\,p(y_t\mid\theta_{k^\ast},\xi_t)\big) .
\]
For Gaussian likelihoods with a linearised forward map
$\mu(\theta,\xi_t)\approx\mu(\theta_{\mathrm{dom}},\xi_t)+J_{\xi_t}(\theta-\theta_{\mathrm{dom}})$,
this reduces to
$\tfrac12(\theta_{\mathrm{dom}}-\theta_{k^\ast})^\top F(\xi_t)(\theta_{\mathrm{dom}}-\theta_{k^\ast})$
with $F(\xi_t)=J_{\xi_t}^\top R^{-1}J_{\xi_t}$, and treating $\theta_{k^\ast}$ as a
posterior draw with covariance $\Sigma_{\mathrm{post}}$ gives the expected criterion
$\tfrac12\,\mathrm{tr}(F(\xi_t)\Sigma_{\mathrm{post}})$, a trace approximation to the
Gaussian-linearised (D-optimal) EIG.

This KL connection holds only for the pairwise case. The general weighted-sum criterion
of Eq.~\eqref{eq:lr-limit} is a sum of likelihood \emph{ratios}, not log-ratios, and
their expectations are constant in the design,
$\mathbb{E}_{y_{t,\mathrm{dom}}}[r_{k\leftarrow\mathrm{dom}}]=\int p(y\mid\theta_k,\xi_t)\,\ud y=1$;
there it is the realised $y_{t,\mathrm{dom}}$, not its expectation, that carries the
design signal.

\section{The dominant particle, and why standard IS fixes do not help}
\label{app:domparticle}
Under a collapsed weight set the dominant particle $\theta_{\mathrm{dom}}$ is the
posterior draw that the NPE most \emph{undercovers}: its self-normalised weight
$\tilde w_j\propto p(\theta_j\mid h)/q_\phi(\theta_j\mid h)$ is largest where the proposal
places too little mass relative to the true posterior, typically a posterior-tail
sample. This is why design selection keeps a signal even as $\mathrm{ESS}\to1$: the
criterion of Eq.~\eqref{eq:lr-limit} discriminates this tail representative from the
bulk rather than collapsing onto the posterior mean.

Two natural alternatives are worse. Dropping the importance weights and scoring with the
raw NPE particles fails because, once the posterior is tight, all $J$ particles cluster
near $\theta_{\mathrm{dom}}$, every ratio $r_{k\leftarrow j}\approx1$, and the score is
flat across designs; the reweighting is what keeps a discriminative contrastive in the
bank. Standard remedies for heavy-tailed importance weights also do not apply here.
Pareto-smoothed importance sampling~\citep{vehtari2024pareto} assumes a finite-variance
tail (Pareto $\hat k<0.7$), whereas here $\hat k$ is large by construction; truncating
the weights trades the (cancelling) variance for a bias that dominates at
$\mathrm{ESS}/J\sim10^{-3}$; and defensive mixtures that fold the prior back into the
proposal reintroduce exactly the prior-contrastive cost $L_{\mathrm{escape}}$ that PACE
removes. PACE instead accepts the collapse and relies on common-mode cancellation, which
requires only that the weights are shared across candidate designs.

\chapter{Supplementary tables}
\label{app:supp}

Table~\ref{tab:ess-full} reports the deployment-time effective sample size across
steps. It declines on every system but does not track the performance gap: see
Section~\ref{sec:r-perstep}. For the per-step increments $\Delta S_t$, the
cumulative $\sum_t\Delta S_t$ for PACE is $+13.06$ on CES (DAD $+12.05$, random
$+9.05$) and $+14.83$ on SEIR R=3 (DAD $+7.46$, random $+14.15$), and exceeds random
on every system in the suite.

\begin{table}[h]
\centering\footnotesize
\caption[Effective sample size trajectories]{Median per-step effective sample size from
production rollouts ($J=1000$; reported at $J=1000$ for ASM1 for comparability). The
decline is monotone from the first step and universal; all eight systems reach
$\mathrm{ESS}\le3$ by the final step, yet the decline does not track the performance
gap (Section~\ref{sec:r-perstep}). Only PK 8-comp (omitted, $\mathrm{ESS}/J\approx12$
to $23\%$ throughout) retains a high effective sample size.}
\label{tab:ess-full}
\begin{tabular}{l c >{\raggedright\arraybackslash}p{8.2cm} l}
\toprule
System & $T$ & ESS per step (median) & $\Delta_{\mathrm{sPCE}}$ \\
\midrule
Monod & 14 & 841, 612, 330, 125, 68, 39, 36, 18, 17, 10, 5, 3, 2, 1 & $-0.66$ vs DAD \\
CES & 10 & 9, 2, 2, 2, 3, 2, 2, 3, 2, 2 & $+0.93$ vs DAD \\
Src.\ loc.\ $p{=}2$ & 10 & 40, 29, 15, 12, 6, 3, 2, 2, 2, 1 & $+1.52$ vs DAD \\
Src.\ loc.\ $p{=}5$ & 30 & 93, 75, 39, 37, 32, 29, 21, 18, 13, 11, 8, 7, 6, 5, 6, 5, 4, 4, 4, 3, 3, 3, 3, 3, 3, 2, 2, 2, 2, 2 & $+1.07$ vs DAD \\
PI3K & 8 & 561, 285, 87, 17, 8, 3, 3, 2 & $-0.14$ vs DAD \\
SEIR R=3 & 10 & 20, 8, 3, 2, 2, 1, 1, 2, 1, 1 & $+7.31$ vs DAD \\
ASM1 & 5 & 77, 16, 6, 3, 3 & $+0.89$ vs rand. \\
Goodwin & 10 & 169, 81, 37, 23, 12, 8, 6, 5, 6, 2 & $-0.48$ vs rand. \\
\bottomrule
\end{tabular}
\end{table}

% AUTO-GENERATED by scripts/_gen_perstep_table.py — do not edit by hand.
% Re-run after adding systems to outputs/_perstep_compiled/.

\section{Per-step EIG contribution under ESS collapse}
\label{sec:perstep-dS}

Table~\ref{tab:ess-full} shows the IS-weight ESS collapsing across BED steps;
Table~\ref{tab:perstep-dS} shows the per-step \emph{information} still
contributed despite that collapse. We report per-step increments
$\Delta S_t = S_t - S_{t-1}$, where $S_t$ is the sPCE bound
[Eq.~\eqref{eq:spce}] evaluated on the length-$t$ prefix $\hist_t$ (with $S_0=0$
and $S_T=\mathrm{sPCE}_T$). By the chain rule of mutual information these
telescope, $\sum_{t=1}^{T}\Delta S_t = S_T$, approximately recovering the trajectory bound
(last column of Table~\ref{tab:perstep-dS}); each $\Delta S_t$ estimates the
conditional EIG $I(\theta;y_t\mid\hist_{t-1})$ [Eq.~\eqref{eq:step-eig}] at the
realized $y_t$. Crucially, $S_t$ weights the contrastive denominator by the
running posterior $w^{(t-1)}_\ell\propto p(\hist_{t-1}\mid\theta_\ell)$, the same
importance weights of \S\ref{sec:m-costtransfer} whose ESS collapses as the
posterior contracts, so $\Delta S_t$ measures per-step information \emph{within}
the degenerate-weight regime rather than from a fresh prior-contrastive
estimate.\footnote{This conditional increment differs from the marginal
per-step sPCE (uniform contrastives, targeting $I(\theta;y_t)$), whose sum
overcounts $\mathrm{sPCE}_T$ by the total correlation among observations and is
blind to the collapse. As a numerical check, $\sum_t\Delta S_t$ matches the
trajectory sPCE of Table~\ref{tab:results} to within MC error on cells with
matching evaluation protocol.}

\textbf{Findings.}
\emph{(i) Per-step EIG decay tracks posterior concentration, not ESS.}
$\Delta S_t$ decays across steps for every method, but this reflects the
information remaining in a contracting posterior, \emph{not} the IS-weight ESS
collapse of Table~\ref{tab:ess-full}, from which it is empirically decoupled.
On CES the ESS pins at ${\approx}2$ from step~2 onward, yet $70\%$ of PACE's total
information is still acquired at $\mathrm{ESS}\leq 5$ (steps~3--5 alone, at
$\mathrm{ESS}\leq 3$, contribute $34\%$); on source-location $p{=}2$ the
PACE$-$random per-step gap holds at $\approx{+}0.3$ nat through the late steps
(8--10) where $\mathrm{ESS}\leq 5$. SEIR~R=3 is the extreme case: its ESS
collapses to $1$ by step~$10$ (Table~\ref{tab:ess-full}), yet the per-step
PACE$-$DAD gap \emph{grows} from $+0.33$ nat (first third) to $+2.02$ nat
(last third), as DAD's late designs collapse (Table~\ref{tab:dad-sweep})
while PACE's stay informative. Conversely, PK 8-comp keeps
$\mathrm{ESS}/J\geq 12\%$ throughout (PI3K, once thought similar, in fact
collapses to $\mathrm{ESS}\approx 2$ under the deployed random NPE), yet its
$\Delta S_t$ decays just as steeply. Per-step design quality follows the information remaining, not the
ESS: the per-step, mechanistic counterpart of the aggregate result that ESS
does not predict performance (\S\ref{sec:r-perstep}).
\emph{(ii) The cumulative advantage holds throughout.} PACE's
$\Sigma_t\Delta S_t$ exceeds random by $+0.5$ to $+5.6$
nat on each system in this table, and exceeds DAD on CES, SEIR~R=3, source-location $p{=}2$ and $p{=}5$,
PK~8-comp, and PK~1-comp $T{=}5$; it ties PK~1-comp $T{=}3$ and trails DAD on
Monod, PI3K, and source-location $p{=}8$ (consistent with
Table~\ref{tab:results}). The few late steps
at which PACE dips below random do so by small margins (at most $0.10$ nat, on
SEIR~R=3 and PK~1-comp $T{=}5$) and
never erode the cumulative gap; PACE never \emph{collapses} below baseline under
ESS collapse, consistent with the common-mode-cancellation guarantee of
\S\ref{sec:m-lowess}.
\emph{(iii) Front-loaded, with saturation-driven late convergence.}
Source-location $p{=}2$ and PI3K retain a clear late-step gap over random
($\geq 0.26$ nat at the final step); CES, Monod, source-location $p{=}5$, and
PK 8-comp converge to $\approx 0$ late as the posterior is near-resolved and
\emph{all} methods stop gaining: task saturation, not PACE degradation. (The
PK 1-comp random baseline, $n{=}45$, is noisier per step.)

This is the per-step, ESS-aware complement to the aggregate result that ESS
does not predict performance (\S\ref{sec:r-perstep}): even as
$\mathrm{ESS}\to 1$, PACE's incremental information stays non-negative and at
least matches the baselines at every step. (ASM1 and Goodwin are omitted
here pending their chain-rule re-evaluation; ASM1's $T{=}5$ horizon gives
little late-step resolution regardless.)

\begin{table}[h]
\centering\scriptsize
\setlength{\tabcolsep}{4pt}\renewcommand{\arraystretch}{0.92}
\caption{Per-step chain-rule EIG contribution $\Delta S_t = S_t - S_{t-1}$ (mean over $n$ rollouts; $S_t$ = sPCE on the length-$t$ prefix). All PACE rows use the deployed PACE-ftk (CEM) configuration of Table~\ref{tab:results} (the same per-rollout runs; PACE-grad is the ablation in Table~\ref{tab:ablation}, not shown here). Read alongside Table~\ref{tab:ess-full} (ESS per step). Last column $\Sigma_t \Delta S_t$ = trajectory sPCE, approximately recovering Table~\ref{tab:results} to within MC error. $\Delta S_t$ decays for all methods (posterior concentration, decoupled from the ESS of Table~\ref{tab:ess-full}). PK 1-comp random uses $n{=}45$ rollouts and is noisier per step. }
\label{tab:perstep-dS}
\begin{tabular}{ll>{\raggedright\tiny\arraybackslash}p{8.4cm}r}
\toprule
System & Method & $\Delta S_t$ per step (mean), $t=1\ldots T$ & $\Sigma_t\Delta S_t$ \\
\midrule
  CES & PACE & 3.93\,$\to$\allowbreak 2.83\,$\to$\allowbreak 2.05\,$\to$\allowbreak 1.50\,$\to$\allowbreak 0.88\,$\to$\allowbreak 0.60\,$\to$\allowbreak 0.35\,$\to$\allowbreak 0.37\,$\to$\allowbreak 0.27\,$\to$\allowbreak 0.28 & +13.06 \\
   & DAD & 3.70\,$\to$\allowbreak 2.98\,$\to$\allowbreak 1.38\,$\to$\allowbreak 0.82\,$\to$\allowbreak 0.89\,$\to$\allowbreak 0.63\,$\to$\allowbreak 0.59\,$\to$\allowbreak 0.47\,$\to$\allowbreak 0.36\,$\to$\allowbreak 0.23 & +12.05 \\
   & random & 2.41\,$\to$\allowbreak 1.76\,$\to$\allowbreak 1.55\,$\to$\allowbreak 0.97\,$\to$\allowbreak 0.57\,$\to$\allowbreak 0.46\,$\to$\allowbreak 0.41\,$\to$\allowbreak 0.39\,$\to$\allowbreak 0.28\,$\to$\allowbreak 0.24 & +9.05 \\
  \cmidrule(l){1-4}
  SEIR R=3 & PACE & 5.47\,$\to$\allowbreak 2.46\,$\to$\allowbreak 1.53\,$\to$\allowbreak 1.40\,$\to$\allowbreak 1.03\,$\to$\allowbreak 0.68\,$\to$\allowbreak 0.78\,$\to$\allowbreak 0.59\,$\to$\allowbreak 0.46\,$\to$\allowbreak 0.42 & +14.83 \\
   & DAD & 4.64\,$\to$\allowbreak 2.11\,$\to$\allowbreak 1.74\,$\to$\allowbreak 1.46\,$\to$\allowbreak 1.03\,$\to$\allowbreak 0.75\,$\to$\allowbreak 0.31\,$\to$\allowbreak -0.29\,$\to$\allowbreak -1.54\,$\to$\allowbreak -2.74 & +7.46 \\
   & random & 4.82\,$\to$\allowbreak 2.17\,$\to$\allowbreak 1.49\,$\to$\allowbreak 1.42\,$\to$\allowbreak 1.13\,$\to$\allowbreak 0.77\,$\to$\allowbreak 0.84\,$\to$\allowbreak 0.60\,$\to$\allowbreak 0.47\,$\to$\allowbreak 0.46 & +14.15 \\
  \cmidrule(l){1-4}
  Monod & PACE & 0.07\,$\to$\allowbreak 0.23\,$\to$\allowbreak 0.44\,$\to$\allowbreak 0.69\,$\to$\allowbreak 0.33\,$\to$\allowbreak 0.38\,$\to$\allowbreak 0.44\,$\to$\allowbreak 0.28\,$\to$\allowbreak 0.38\,$\to$\allowbreak 0.40\,$\to$\allowbreak 0.51\,$\to$\allowbreak 0.47\,$\to$\allowbreak 0.39\,$\to$\allowbreak 0.29 & +5.30 \\
   & DAD & 0.03\,$\to$\allowbreak 0.16\,$\to$\allowbreak 0.42\,$\to$\allowbreak 0.47\,$\to$\allowbreak 0.61\,$\to$\allowbreak 0.49\,$\to$\allowbreak 0.58\,$\to$\allowbreak 0.51\,$\to$\allowbreak 0.49\,$\to$\allowbreak 0.48\,$\to$\allowbreak 0.49\,$\to$\allowbreak 0.48\,$\to$\allowbreak 0.37\,$\to$\allowbreak 0.32 & +5.88 \\
   & random & 0.04\,$\to$\allowbreak 0.12\,$\to$\allowbreak 0.23\,$\to$\allowbreak 0.47\,$\to$\allowbreak 0.27\,$\to$\allowbreak 0.34\,$\to$\allowbreak 0.48\,$\to$\allowbreak 0.33\,$\to$\allowbreak 0.41\,$\to$\allowbreak 0.43\,$\to$\allowbreak 0.53\,$\to$\allowbreak 0.42\,$\to$\allowbreak 0.36\,$\to$\allowbreak 0.32 & +4.76 \\
  \cmidrule(l){1-4}
  PI3K & PACE & 2.11\,$\to$\allowbreak 0.99\,$\to$\allowbreak 0.80\,$\to$\allowbreak 0.68\,$\to$\allowbreak 0.49\,$\to$\allowbreak 0.38\,$\to$\allowbreak 0.44\,$\to$\allowbreak 0.48 & +6.36 \\
   & DAD & 1.77\,$\to$\allowbreak 1.23\,$\to$\allowbreak 0.91\,$\to$\allowbreak 0.77\,$\to$\allowbreak 0.55\,$\to$\allowbreak 0.50\,$\to$\allowbreak 0.46\,$\to$\allowbreak 0.34 & +6.53 \\
   & random & 1.46\,$\to$\allowbreak 0.17\,$\to$\allowbreak 0.08\,$\to$\allowbreak 0.13\,$\to$\allowbreak 0.16\,$\to$\allowbreak 0.13\,$\to$\allowbreak 0.11\,$\to$\allowbreak 0.20 & +2.45 \\
  \cmidrule(l){1-4}
  PK 1-comp $T{=}3$ & PACE & 1.19\,$\to$\allowbreak 0.95\,$\to$\allowbreak 0.48 & +2.62 \\
   & DAD & 1.01\,$\to$\allowbreak 0.91\,$\to$\allowbreak 0.71 & +2.64 \\
   & random & 1.26\,$\to$\allowbreak 0.51\,$\to$\allowbreak 0.29 & +2.05 \\
  \cmidrule(l){1-4}
  PK 1-comp $T{=}5$ & PACE & 1.20\,$\to$\allowbreak 0.95\,$\to$\allowbreak 0.47\,$\to$\allowbreak 0.34\,$\to$\allowbreak 0.32 & +3.27 \\
   & DAD & 1.23\,$\to$\allowbreak 0.33\,$\to$\allowbreak 0.49\,$\to$\allowbreak 0.81\,$\to$\allowbreak 0.33 & +3.19 \\
   & random & 1.26\,$\to$\allowbreak 0.51\,$\to$\allowbreak 0.29\,$\to$\allowbreak 0.42\,$\to$\allowbreak 0.23 & +2.70 \\
  \cmidrule(l){1-4}
  PK 8-comp $T{=}5$ & PACE & 1.12\,$\to$\allowbreak 0.48\,$\to$\allowbreak 0.31\,$\to$\allowbreak 0.19\,$\to$\allowbreak 0.16 & +2.26 \\
   & DAD & 0.42\,$\to$\allowbreak 0.64\,$\to$\allowbreak 0.53\,$\to$\allowbreak 0.30\,$\to$\allowbreak 0.22 & +2.11 \\
   & random & 0.35\,$\to$\allowbreak 0.19\,$\to$\allowbreak 0.15\,$\to$\allowbreak 0.17\,$\to$\allowbreak 0.11 & +0.96 \\
  \cmidrule(l){1-4}
  Src.\ loc.\ $p{=}2$ & PACE & 1.11\,$\to$\allowbreak 1.14\,$\to$\allowbreak 1.04\,$\to$\allowbreak 1.16\,$\to$\allowbreak 0.97\,$\to$\allowbreak 0.91\,$\to$\allowbreak 0.78\,$\to$\allowbreak 0.75\,$\to$\allowbreak 0.65\,$\to$\allowbreak 0.58 & +9.09 \\
   & DAD & 0.94\,$\to$\allowbreak 0.94\,$\to$\allowbreak 0.95\,$\to$\allowbreak 0.87\,$\to$\allowbreak 0.77\,$\to$\allowbreak 0.75\,$\to$\allowbreak 0.76\,$\to$\allowbreak 0.71\,$\to$\allowbreak 0.54\,$\to$\allowbreak 0.37 & +7.61 \\
   & random & 0.70\,$\to$\allowbreak 0.67\,$\to$\allowbreak 0.50\,$\to$\allowbreak 0.42\,$\to$\allowbreak 0.50\,$\to$\allowbreak 0.46\,$\to$\allowbreak 0.39\,$\to$\allowbreak 0.39\,$\to$\allowbreak 0.27\,$\to$\allowbreak 0.32 & +4.62 \\
  \cmidrule(l){1-4}
  Src.\ loc.\ $p{=}5$ & PACE & 0.55\,$\to$\allowbreak 0.57\,$\to$\allowbreak 0.35\,$\to$\allowbreak 0.48\,$\to$\allowbreak 0.56\,$\to$\allowbreak 0.48\,$\to$\allowbreak 0.42\,$\to$\allowbreak 0.53\,$\to$\allowbreak 0.34\,$\to$\allowbreak 0.35\,$\to$\allowbreak 0.48\,$\to$\allowbreak 0.28\,$\to$\allowbreak 0.34\,$\to$\allowbreak 0.36\,$\to$\allowbreak 0.36\,$\to$\allowbreak 0.33\,$\to$\allowbreak 0.29\,$\to$\allowbreak 0.31\,$\to$\allowbreak 0.27\,$\to$\allowbreak 0.32\,$\to$\allowbreak 0.27\,$\to$\allowbreak 0.26\,$\to$\allowbreak 0.22\,$\to$\allowbreak 0.20\,$\to$\allowbreak 0.21\,$\to$\allowbreak 0.16\,$\to$\allowbreak 0.18\,$\to$\allowbreak 0.10\,$\to$\allowbreak 0.14\,$\to$\allowbreak 0.12 & +9.82 \\
   & DAD & 0.46\,$\to$\allowbreak 0.34\,$\to$\allowbreak 0.36\,$\to$\allowbreak 0.41\,$\to$\allowbreak 0.44\,$\to$\allowbreak 0.38\,$\to$\allowbreak 0.32\,$\to$\allowbreak 0.35\,$\to$\allowbreak 0.35\,$\to$\allowbreak 0.27\,$\to$\allowbreak 0.43\,$\to$\allowbreak 0.44\,$\to$\allowbreak 0.34\,$\to$\allowbreak 0.29\,$\to$\allowbreak 0.33\,$\to$\allowbreak 0.35\,$\to$\allowbreak 0.26\,$\to$\allowbreak 0.27\,$\to$\allowbreak 0.25\,$\to$\allowbreak 0.28\,$\to$\allowbreak 0.22\,$\to$\allowbreak 0.22\,$\to$\allowbreak 0.22\,$\to$\allowbreak 0.17\,$\to$\allowbreak 0.17\,$\to$\allowbreak 0.21\,$\to$\allowbreak 0.16\,$\to$\allowbreak 0.21\,$\to$\allowbreak 0.14\,$\to$\allowbreak 0.14 & +8.79 \\
   & random & 0.25\,$\to$\allowbreak 0.22\,$\to$\allowbreak 0.15\,$\to$\allowbreak 0.17\,$\to$\allowbreak 0.17\,$\to$\allowbreak 0.15\,$\to$\allowbreak 0.16\,$\to$\allowbreak 0.10\,$\to$\allowbreak 0.12\,$\to$\allowbreak 0.21\,$\to$\allowbreak 0.13\,$\to$\allowbreak 0.12\,$\to$\allowbreak 0.11\,$\to$\allowbreak 0.07\,$\to$\allowbreak 0.17\,$\to$\allowbreak 0.13\,$\to$\allowbreak 0.21\,$\to$\allowbreak 0.16\,$\to$\allowbreak 0.17\,$\to$\allowbreak 0.17\,$\to$\allowbreak 0.17\,$\to$\allowbreak 0.23\,$\to$\allowbreak 0.08\,$\to$\allowbreak 0.02\,$\to$\allowbreak 0.11\,$\to$\allowbreak 0.04\,$\to$\allowbreak 0.10\,$\to$\allowbreak 0.12\,$\to$\allowbreak 0.10\,$\to$\allowbreak 0.10 & +4.20 \\
  \cmidrule(l){1-4}
  Src.\ loc.\ $p{=}8$ & PACE & 0.35\,$\to$\allowbreak 0.28\,$\to$\allowbreak 0.18\,$\to$\allowbreak 0.30\,$\to$\allowbreak 0.23\,$\to$\allowbreak 0.16\,$\to$\allowbreak 0.18\,$\to$\allowbreak 0.19\,$\to$\allowbreak 0.21\,$\to$\allowbreak 0.17\,$\to$\allowbreak 0.18\,$\to$\allowbreak 0.24\,$\to$\allowbreak 0.18\,$\to$\allowbreak 0.17\,$\to$\allowbreak 0.21\,$\to$\allowbreak 0.19\,$\to$\allowbreak 0.16\,$\to$\allowbreak 0.18\,$\to$\allowbreak 0.20\,$\to$\allowbreak 0.11\,$\to$\allowbreak 0.17\,$\to$\allowbreak 0.26\,$\to$\allowbreak 0.17\,$\to$\allowbreak 0.15\,$\to$\allowbreak 0.18\,$\to$\allowbreak 0.09\,$\to$\allowbreak 0.12\,$\to$\allowbreak 0.16\,$\to$\allowbreak 0.14\,$\to$\allowbreak 0.16 & +5.66 \\
   & DAD & 0.30\,$\to$\allowbreak 0.21\,$\to$\allowbreak 0.19\,$\to$\allowbreak 0.29\,$\to$\allowbreak 0.23\,$\to$\allowbreak 0.23\,$\to$\allowbreak 0.20\,$\to$\allowbreak 0.17\,$\to$\allowbreak 0.21\,$\to$\allowbreak 0.22\,$\to$\allowbreak 0.23\,$\to$\allowbreak 0.25\,$\to$\allowbreak 0.22\,$\to$\allowbreak 0.26\,$\to$\allowbreak 0.22\,$\to$\allowbreak 0.22\,$\to$\allowbreak 0.21\,$\to$\allowbreak 0.20\,$\to$\allowbreak 0.17\,$\to$\allowbreak 0.23\,$\to$\allowbreak 0.16\,$\to$\allowbreak 0.16\,$\to$\allowbreak 0.16\,$\to$\allowbreak 0.14\,$\to$\allowbreak 0.19\,$\to$\allowbreak 0.17\,$\to$\allowbreak 0.15\,$\to$\allowbreak 0.18\,$\to$\allowbreak 0.13\,$\to$\allowbreak 0.12 & +6.00 \\
   & random & 0.09\,$\to$\allowbreak 0.07\,$\to$\allowbreak 0.02\,$\to$\allowbreak 0.09\,$\to$\allowbreak 0.08\,$\to$\allowbreak 0.04\,$\to$\allowbreak 0.08\,$\to$\allowbreak 0.07\,$\to$\allowbreak -0.01\,$\to$\allowbreak 0.09\,$\to$\allowbreak 0.04\,$\to$\allowbreak 0.04\,$\to$\allowbreak 0.04\,$\to$\allowbreak 0.05\,$\to$\allowbreak 0.08\,$\to$\allowbreak 0.08\,$\to$\allowbreak 0.05\,$\to$\allowbreak 0.10\,$\to$\allowbreak 0.06\,$\to$\allowbreak 0.14\,$\to$\allowbreak 0.08\,$\to$\allowbreak 0.10\,$\to$\allowbreak 0.07\,$\to$\allowbreak 0.03\,$\to$\allowbreak 0.03\,$\to$\allowbreak 0.04\,$\to$\allowbreak 0.05\,$\to$\allowbreak 0.05\,$\to$\allowbreak 0.09\,$\to$\allowbreak 0.06 & +1.87 \\
  \bottomrule
\end{tabular}
\end{table}


\chapter{Diagnostic protocols}
\label{app:diagnostics}

\paragraph{Bootstrap pick-stability.}
On ASM1, for several steps of a ten-step rollout, draw $J=500$ particles,
bootstrap-resample $B=200$ times, recompute the PACE scores for $50$ candidates, and
record the top-1 agreement. The median is $45\%$, about $23\times$ the chance rate
of $1/50$. Residual variation comes from resampling changing which particle
dominates.

\paragraph{ESS-regime interpolation.}
Classifying each (system, step) cell by $\mathrm{ESS}/J$ as EIG-estimation
($\ge10\%$), transition ($1$--$10\%$), or dominant-particle ($<1\%$),
PACE spans a more than hundredfold range across systems without any user intervention,
and individual systems traverse the regimes within a single rollout as the posterior
tightens (Table~\ref{tab:ess-regime}). PK 8-comp, with tight priors, is the only system
that stays in the EIG-estimation regime throughout; PI3K traverses all three;
source-location and ASM1 reach the dominant-particle regime by their final steps yet
still deliver positive design value. This is the smooth interpolation of
Section~\ref{sec:m-cancellation} made empirical.

\begin{table}[h]
\centering\small
\caption[ESS regimes]{Range of $\mathrm{ESS}/J$ across steps and the regime traversed,
classified as EIG-estimation ($\ge10\%$), transition ($1$--$10\%$), or
dominant-particle ($<1\%$). The regime is set by system informativeness (the
prior-to-posterior volume ratio), not parameter dimension.}
\label{tab:ess-regime}
\begin{tabular}{lccl}
\toprule
System & $d_\theta$ & $\mathrm{ESS}/J$ range & Regime traversed \\
\midrule
PK 8-comp & 20 & 12--23\% & EIG estimation throughout \\
PI3K & 6 & 0.2--56\% & EIG $\to$ transition $\to$ dominant-particle \\
Src.\ loc.\ $p{=}2$ & 4 & 0.1--4\% & transition $\to$ dominant-particle \\
ASM1 & 12 & 0.3--8\% & transition $\to$ dominant-particle \\
\bottomrule
\end{tabular}
\end{table}

\paragraph{Raw versus resampled coverage.}
Computed in Section~\ref{sec:r-diag}, Table~\ref{tab:coverage}: raw NPE coverage of
the truth stays in $75$ to $90\%$ while resampled coverage collapses at late steps.

\paragraph{Per-step calibration.}
Table~\ref{tab:sbc} reports simulation-based calibration KS distances on raw NPE
samples and on resampled samples. Raw KS is flat or improving across steps; the
resampled KS degrades, confirming the resampling artifact rather than a proposal
failure. Simulation-based calibration tests average-case calibration; the raw
coverage measured at deployment histories is the complementary conditional check.

\begin{table}[h]
\centering\small
\caption[Per-step calibration]{Simulation-based calibration KS distance (smaller is
better) on raw NPE samples (tests $q_\phi$ directly) versus importance-resampled
samples. Upper block: stiff-ODE systems with severe ESS collapse ($n{=}20$ synthetic
trajectories, $J{=}1000$). Lower block: unsaturated systems from controlled diagnostics
(IS KS not computed). Raw KS stays low and stable across steps; IS KS degrades as the
effective sample size collapses. $^\ddagger$The source-location $p{=}8$ diagnostic uses
a $15$-step horizon; the production system uses $T{=}30$.}
\label{tab:sbc}
\begin{tabular}{lcccc}
\toprule
System & step & ESS & raw KS & IS KS \\
\midrule
CES ($d_\theta{=}4$, $T{=}10$)
  & 1 & $65.8$ & $0.17$ & $0.20$ \\
  & 5 & $6.7$ & $0.14$ & $0.33$ \\
  & 10 & $10.7$ & $0.13$ & $0.39$ \\
ASM1 ($d_\theta{=}12$, $T{=}5$)
  & 1 & $97.6$ & $0.18$ & $0.18$ \\
  & 3 & $14.1$ & $0.17$ & $0.20$ \\
  & 5 & $8.4$ & $0.17$ & $0.24$ \\
\midrule
PI3K ($d_\theta{=}6$, $T{=}8$)
  & 1 & -- & $0.14$ & -- \\
  & 8 & -- & $0.13$ & -- \\
PK 8-comp ($d_\theta{=}20$, $T{=}5$)
  & 1 & -- & $0.16$ & -- \\
  & 5 & -- & $0.15$ & -- \\
Src.\ loc.\ $p{=}8$ ($d_\theta{=}16$)$^\ddagger$
  & 1 & -- & $0.17$ & -- \\
  & 15 & -- & $0.16$ & -- \\
\bottomrule
\end{tabular}
\end{table}

\chapter{Wall-clock cost}
\label{app:walltime}

Table~\ref{tab:walltime} compares the one-time training cost and the per-step
deployment cost. PACE trains its posterior estimator once in $3.6$ to $12.8$ minutes
and then selects each design in under four seconds on eleven of twelve systems (ASM1,
a stiff $13$-state ODE, is the exception at $40.5$\,s). Policy training ranges from
minutes to sixteen hours and is infeasible on the stiffest systems.

\paragraph{Hardware.} All Snellius runs use a single NVIDIA A100-SXM4-40GB on the
\texttt{gpu\_a100} partition; this covers every PACE per-step measurement, all DAD
training, and all NPE training except ASM1 and source-location, whose NPEs were trained
on a local desktop GPU (marked $\dagger$; an A100 would be faster, which only
strengthens the comparison). Drawing $J{=}1000$ particles from a trained NPE takes
about $11$\,ms, negligible against the ODE-solve cost.

\begin{table}[h]
\centering\small
\caption[Wall-clock cost]{Per-step design time, one-time NPE training, and one-time
policy training, with the experiment timescale for context. A dagger marks a local-GPU
measurement; an asterisk marks unstable DAD training (CES diverged to NaN before
convergence; the reported time is a partial run).}
\label{tab:walltime}
\begin{tabular}{lrrrll}
\toprule
System & $T$ & PACE/step & NPE train & DAD train & Expt.\ timescale \\
\midrule
CES & 10 & $0.09$\,s & $3.7$\,min & $9$\,min$^*$ & minutes \\
ASM1 & 5 & $40.5$\,s & $12.8$\,min$^\dagger$ & infeasible & hours \\
SEIR R=3 & 10 & ${\sim}1$\,s & $3.6$\,min & $9$\,min (useless) & days--weeks \\
Monod & 14 & $0.22$\,s & $7.2$\,min & $16$\,hr & hours \\
Goodwin & 10 & $0.16$\,s & $5.1$\,min & failed & -- \\
Src.\ loc.\ $p{=}2$ & 10 & $0.13$\,s & $3$--$12$\,min$^\dagger$ & $20$\,min & -- \\
Src.\ loc.\ $p{=}5$ & 30 & $0.15$\,s & $3$--$12$\,min$^\dagger$ & $1.5$\,hr & -- \\
Src.\ loc.\ $p{=}8$ & 30 & $0.16$\,s & $3$--$12$\,min$^\dagger$ & $1.5$\,hr & -- \\
PK 1-comp & 5 & $1.21$\,s & $8.5$--$9$\,min & $7$\,min & hours \\
PK 8-comp & 5 & $0.34$\,s & $9.6$--$11$\,min & $31$\,min & hours \\
PI3K & 8 & $3.34$\,s & $8$--$9$\,min & $3.2$\,hr & hours \\
\bottomrule
\end{tabular}
\end{table}

\chapter{System specifications}
\label{app:systems}

Every system used in the thesis is defined here, one section each, in a common format:
the forward model and its constants, the prior, and the observation and design spaces.
Forward models, parameters, priors, and noise are taken from the implementation;
per-method experiment configurations (horizon $T$, particle count $J$, candidates
$N_{\mathrm{cand}}$, CEM rounds $C$, evaluation budget $L_{\mathrm{eval}}$) are collected
in Table~\ref{tab:config}. Dynamic systems are integrated with an explicit (Tsit5)
solver, except where stiffness is noted. Throughout, $\mathcal N$ denotes a normal
distribution and $\mathrm{LogNormal}(\log m,\sigma^2)$ a variable whose natural log is
$\mathcal N(\log m,\sigma^2)$.

\section{CES (constant elasticity of substitution)}
\textbf{Forward model.} A static utility-comparison model~\citep{foster2021deep,blau2022optimizing}.
A design $\xi=(\xi^{(1)},\xi^{(2)})\in[0.01,100]^3\times[0.01,100]^3$ presents two
commodity bundles. With parameters $\theta=(\rho,\alpha_1,\alpha_2,\log u)$ the
CES utility of a bundle $x$ is $U(x)=\big(\sum_{i=1}^3\alpha_i x_i^{\rho}\big)^{1/\rho}$,
and the response is the censored sigmoid-normal
\[
y=\sigma(\eta),\quad \eta\sim\mathcal N(\mu_\eta,\sigma_\eta^2),\quad
\mu_\eta=u\big(U(\xi^{(1)})-U(\xi^{(2)})\big),\quad
\sigma_\eta=u\,\sigma_0\big(1+\|\xi^{(1)}-\xi^{(2)}\|_2\big),
\]
with base noise $\sigma_0=0.005$ and $y$ censored to $[\epsilon,1-\epsilon]$,
$\epsilon=2^{-22}$.
\textbf{Prior.} $\rho\sim\mathrm{Beta}(1,1)=\mathrm{Uniform}(0,1)$ for elasticity,
$(\alpha_1,\alpha_2,\alpha_3)\sim\mathrm{Dirichlet}(1,1,1)$ for the simplex weights
($\alpha_3=1-\alpha_1-\alpha_2$), and $\log u\sim\mathcal N(1,3^2)$ for the scale;
$d_\theta=4$.
\textbf{Observation and design.} $d_y=1$, $d_\xi=6$. The censored boundary mass at
$0$ and $1$ is what destabilises policy training on this system.

\section{Source location}
\textbf{Forward model.} $K=2$ point sources at locations $\theta=(\theta_1,\dots,\theta_K)$,
$\theta_k\in\mathbb{R}^p$~\citep{foster2021deep}. A sensor at design $\xi\in[-3,3]^p$
reads the log total signal
\[
G(\xi,\theta)=\log\!\Big(\gamma+\sum_{k=1}^{K}\frac{1}{\sigma_s^2+\|\xi-\theta_k\|_2^2}\Big),
\qquad y=G(\xi,\theta)+\varepsilon,\quad \varepsilon\sim\mathcal N(0,0.5^2),
\]
with baseline $\gamma=0.1$ and source bandwidth $\sigma_s^2=10^{-4}$. A multi-resolution
variant returns $d_y=3$ channels at bandwidths $\sigma_s^2\in\{0.001,0.1,1.0\}$.
\textbf{Prior.} $\theta_k\sim\mathcal N(0,I_p)$, with the sources ordered by radial
distance $\|\theta_1\|\le\dots\le\|\theta_K\|$ to break label-switching symmetry; this
gives $d_\theta=Kp=2p$.
\textbf{Observation and design.} $d_\xi=p$, with $p\in\{2,5,8\}$ (so $d_\theta\in\{4,10,16\}$).

\section{Pharmacokinetics: one compartment}
\textbf{Forward model.} The iDAD one-compartment model~\citep{ivanova2021implicit}. With
$\theta=(k_a,k_e,V)$ (absorption rate, elimination rate, volume of distribution) the
plasma concentration at sampling time $\xi\in[0,24]$\,h after a dose $D=400$\,mg is the
closed form
\[
C(\xi,\theta)=\frac{D}{V}\,\frac{k_a}{k_a-k_e}\big(e^{-k_e\xi}-e^{-k_a\xi}\big),
\qquad y\sim\mathcal N\!\big(C,\,\sigma^2(\xi,\theta)\big),\quad
\sigma^2=(0.1\,C)^2+0.316^2,
\]
a heteroscedastic noise with $10\%$ relative and $\sqrt{0.1}$ absolute components.
\textbf{Prior.} $\log\theta\sim\mathcal N\big((\log 1.0,\log 0.1,\log 20.0),\,0.05\,I_3\big)$.
\textbf{Observation and design.} $d_y=1$, $d_\xi=1$. The forward map exposes three
kinetic parameters; the results tables report $d_\theta=4$, a discrepancy to reconcile
(it counts an additional inferred scale).% TODO(verify PK 1-comp parameter count vs tab:config)

\section{Pharmacokinetics: eight-compartment PBPK}
\textbf{Forward model.} A permeability-limited physiologically based model with seven
tissues, central blood, and an oral depot ($16$ internal states). The state vector is
$(M_{\mathrm{depot}},C_{\mathrm{blood}},\{C_{\mathrm{vas},i},C_{\mathrm{tis},i}\}_{i=1}^7)$,
governed by the linear system
\begin{align*}
\dot M_{\mathrm{depot}} &= -k_a M_{\mathrm{depot}},\\
\dot C_{\mathrm{blood}} &= \big(\textstyle\sum_i Q_i(C_{\mathrm{vas},i}-C_{\mathrm{blood}})+k_a F M_{\mathrm{depot}}-\mathrm{CL}_r f_{up}C_{\mathrm{blood}}\big)/V_{\mathrm{blood}},\\
\dot C_{\mathrm{vas},i} &= \big(Q_i(C_{\mathrm{blood}}-C_{\mathrm{vas},i})-\mathrm{PS}_i(f_{up}C_{\mathrm{vas},i}-C_{\mathrm{tis},i}/K_{p,i})\big)/V_{\mathrm{vas},i},\\
\dot C_{\mathrm{tis},i} &= \big(\mathrm{PS}_i(f_{up}C_{\mathrm{vas},i}-C_{\mathrm{tis},i}/K_{p,i})-(\mathrm{CL}_h\delta_{\mathrm{liver}}+\mathrm{CL}_{\mathrm{gut}}\delta_{\mathrm{gut}})\,C_{\mathrm{tis},i}/K_{p,i}\big)/V_{\mathrm{tis},i},
\end{align*}
which is solved in closed form by the matrix exponential $y(t)=\exp(A(\theta)t)\,y_0$.
Physiological constants are fixed at $70$-kg adult values (blood $V_{\mathrm{blood}}=5.6$\,L;
tissue volumes, vascular fractions, and flows $Q_i$ summing to a $390$\,L/h cardiac
output). Dosing is a $100$\,mg IV bolus to blood plus a $100$\,mg oral dose to the depot.
\textbf{Prior.} The $d_\theta=20$ parameters are seven partition coefficients
$K_{p,i}\sim\mathrm{LogNormal}(\log 2,0.3^2)$, seven permeabilities
$\mathrm{PS}_i\sim\mathrm{LogNormal}(\log 50,0.5^2)$, hepatic/renal/gut clearances
$\mathrm{CL}_h\sim\mathrm{LogNormal}(\log 20,0.4^2)$,
$\mathrm{CL}_r\sim\mathrm{LogNormal}(\log 10,0.5^2)$,
$\mathrm{CL}_{\mathrm{gut}}\sim\mathrm{LogNormal}(\log 5,0.5^2)$, absorption
$k_a\sim\mathrm{LogNormal}(\log 1.5,0.4^2)$, and bioavailability $F$ and unbound
fraction $f_{up}$, each $\mathrm{sigmoid}(\mathcal N(0,1^2))$. These calibrated, tight
priors give the system its small escape threshold despite $d_\theta=20$.
\textbf{Observation and design.} $d_y=1$ (plasma $C_{\mathrm{blood}}$) at sampling time
$\xi\in[0,24]$\,h, with heteroscedastic noise $\sigma^2=(0.1\,C)^2+0.1$.

\section{PI3K signalling}
\textbf{Forward model.} A four-state model of iRap-induced PIP3
signalling~\citep{liang2024structural} with states
$(y_1,y_2,y_3,y_4)=(\text{irap\_p85},\text{lyn\_irap\_p85},\text{pip3},\text{yfp\_ph\_mem})$
and conservation laws $\text{cfp\_fkbp\_p85}=1-y_1-y_2$, $\text{lyn\_frb}=1-y_2$,
$\text{yfp\_ph\_cyt}=1-y_4$. With designs $u_{\mathrm{iRap}},u_{\mathrm{LY29}}$ the fluxes are
$v_{\mathrm{dimer}}=k_r\,\text{cfp\_fkbp\_p85}\,u_{\mathrm{iRap}}$,
$v_{\mathrm{trimer}}=k_t\,\text{lyn\_frb}\,y_1$,
$v_{\mathrm{pi3k}}=y_2/(1+u_{\mathrm{LY29}}/k_{ly})$, $v_{\mathrm{deg}}=\phi\,y_3$,
$v_{\mathrm{ph}}=k_a\,\text{yfp\_ph\_cyt}\,y_3-\alpha\,y_4$, and
$\dot y_1=v_{\mathrm{dimer}}-v_{\mathrm{trimer}}$, $\dot y_2=v_{\mathrm{trimer}}$,
$\dot y_3=v_{\mathrm{pi3k}}-v_{\mathrm{deg}}$, $\dot y_4=v_{\mathrm{ph}}$, from $y(0)=0$.
\textbf{Prior.} The six parameters $\theta=(k_r,k_t,k_{ly},\phi,k_a,\alpha)$ are
log-normal about the best-fit values $(1.56,0.03,0.16,1.80,0.58,0.11)$ with log-scale
$\sigma=0.5$.
\textbf{Observation and design.} $d_y=2$, the cytosolic readouts
$(\text{cfp\_fkbp\_p85},\text{yfp\_ph\_cyt})$, with $\mathcal N(0,0.02^2)$ noise; designs
$u_{\mathrm{iRap}}\in[0,1]$ and $u_{\mathrm{LY29}}\in[0,25]$.

\section{Monod and Haldane bioreactors}
\textbf{Forward model.} Fed-batch bioreactors~\citep{monod1949growth,strouwen2026deep}
with state $(C_s,C_x,V)$ (substrate, biomass, volume), feed rate $\xi=Q_{\mathrm{in}}\in[0,6]$,
yield $Y_{xs}=0.777$, feed concentration $C_{s,\mathrm{in}}=50$, and
$\dot C_s=-\mu C_x/Y_{xs}+Q_{\mathrm{in}}(C_{s,\mathrm{in}}-C_s)/V$,
$\dot C_x=\mu C_x-Q_{\mathrm{in}}C_x/V$, $\dot V=Q_{\mathrm{in}}$, from $(3,0.3,7)$.
Monod growth is $\mu=\mu_{\max}C_s/(K_s+C_s)$; the Haldane variant adds substrate
inhibition, $\mu=\mu_{\max}C_s/(K_s+C_s+\alpha C_s^2)$.
\textbf{Prior.} Monod ($d_\theta=2$): $\mu_{\max}\sim\mathrm{Uniform}(0.3,0.5)$,
$K_s\sim\mathrm{Uniform}(0.3,0.6)$. Haldane ($d_\theta=1$): $\alpha\sim\mathrm{Uniform}(0,0.15)$
with $\mu_{\max}=0.4$, $K_s=0.45$ fixed.
\textbf{Observation and design.} $d_y=1$ (substrate $C_s$), $\mathcal N(0,0.1^2)$ noise;
$d_\xi=1$. Haldane appears only in the saturation survey.

\section{DC motor}
\textbf{Forward model.} A two-state electromechanical system (current $i$, angular
velocity $\omega$)~\citep{strouwen2026deep}, used only in the saturation survey:
$\dot i=(V_{\mathrm{in}}-R\,i-k\,\omega)/L_a$, $\dot\omega=(k\,i-f\,\omega)/J_m$, from
$(0,0)$, with $R=0.5$, $L_a=0.0045$, friction $f=0.0125$ fixed and voltage design
$V_{\mathrm{in}}\in[0,10]$.
\textbf{Prior.} $d_\theta=2$: $k\sim\mathrm{Uniform}(0.3,0.7)$,
$J_m\sim\mathrm{Uniform}(0.01,0.04)$.
\textbf{Observation.} $d_y=1$ ($\omega$), $\mathcal N(0,1^2)$ noise.

\section{Goodwin oscillator}
\textbf{Forward model.} A three-state negative-feedback oscillator~\citep{strouwen2026deep}
with mRNA, protein, and product $(X,Y,Z)$ and a high-cooperativity Hill repression,
$\dot X=V_s K^n/(K^n+Z^n)-k_d X+u$, $\dot Y=k_t X-k_d Y$, $\dot Z=k_t Y-k_d Z$, from
$(0.1,0.1,0.1)$, with transcriptional input design $u\in[0,1]$.
\textbf{Prior.} $d_\theta=5$: $V_s,K,k_d,k_t$ log-normal about $(0.5,1.0,0.1,1.0)$ with
log-scale $\sigma=0.4$, and the Hill exponent $n\sim\mathcal N(12,1.5^2)$ truncated to
$n>7$. The high $n$ creates the identifiability plateau that makes Goodwin a predicted
counter-case.
\textbf{Observation and design.} $d_y=1$ ($X$), $\mathcal N(0,0.05^2)$ noise; $d_\xi=1$.

\section{Repressilator}
\textbf{Forward model.} A six-state synthetic oscillator (three mRNAs $m_{1:3}$, three
proteins $P_{1:3}$) in cyclic repression, used only in the saturation survey:
$\dot m_i=\gamma-\delta m_i+\alpha K^n/(P_{j(i)}^n+K^n)+u\,[i{=}1]$ and
$\dot P_i=\beta m_i-\mu P_i$, with $P_{j(1)}=P_3,P_{j(2)}=P_1,P_{j(3)}=P_2$, fixed mRNA
decay $\delta=0.0025$, and inducer design $u\in[0,0.5]$ on gene~1.
\textbf{Prior.} $d_\theta=6$: $\alpha,K,\beta,\gamma,\mu$ log-normal about
$(0.5,40,0.05,0.005,0.005)$, and $n\sim\mathcal N(2,0.5^2)$ truncated to $n>1.5$.
\textbf{Observation.} $d_y=1$ ($P_1$), $\mathcal N(0,5^2)$ noise.

\section{ASM1 wastewater}
\textbf{Forward model.} Activated Sludge Model No.\ 1~\citep{henze2000activated}, a
stiff $13$-state ODE solved with an implicit integrator. The state is
$(S_I,S_S,X_I,X_S,X_{BH},X_{BA},X_P,S_O,S_{NO},S_{NH},S_{ND},X_{ND},S_{ALK})$, driven by
eight processes
\begin{align*}
\rho_0 &= \mu_H\tfrac{S_S}{K_S+S_S}\tfrac{S_O}{K_{OH}+S_O}X_{BH}, &
\rho_1 &= \mu_H\tfrac{S_S}{K_S+S_S}\tfrac{K_{OH}}{K_{OH}+S_O}\tfrac{S_{NO}}{K_{NO}+S_{NO}}\eta_g X_{BH},\\
\rho_2 &= \mu_A\tfrac{S_{NH}}{K_{NH}+S_{NH}}\tfrac{S_O}{K_{OA}+S_O}X_{BA}, &
\rho_3 &= b_H X_{BH},\quad \rho_4=b_A X_{BA},\quad \rho_5=k_a S_{ND}X_{BH},\\
\rho_6 &= k_h\tfrac{X_S/X_{BH}}{K_X+X_S/X_{BH}}\big(\tfrac{S_O}{K_{OH}+S_O}+\eta_h\tfrac{K_{OH}}{K_{OH}+S_O}\tfrac{S_{NO}}{K_{NO}+S_{NO}}\big)X_{BH}, &
\rho_7 &= \rho_6\,X_{ND}/X_S,
\end{align*}
combined through the standard ASM1 Petersen stoichiometry into the $13$ state
derivatives (with a continuous-flow term $Q/V=4\,\mathrm{d}^{-1}$, oxygen transfer
$k_La(S_{O,\mathrm{sat}}-S_O)$, $S_{O,\mathrm{sat}}=8$, and waste term $Q_w$ on the
particulates). Fixed coefficients are $Y_A=0.24$, $F_P=0.08$, $i_{XB}=0.086$,
$i_{XP}=0.06$, $K_{NO}=0.5$, $K_{OA}=0.4$, $\eta_h=0.4$, with fixed influent
concentrations.
\textbf{Prior.} The $d_\theta=12$ kinetic parameters are log-normal about their BSM1
nominal values: $\mu_H\,(6.0)$, $K_S\,(20)$, $K_{OH}\,(0.2)$, $\mu_A\,(0.8)$,
$K_{NH}\,(1.0)$, $b_H\,(0.62)$, $b_A\,(0.15)$, $\eta_g\,(0.8)$, $k_h\,(3.0)$,
$K_X\,(0.03)$, $k_a\,(0.08)$, and yield $Y_H\,(0.67)$, with log-scale $\sigma$ between
$0.19$ and $0.62$ per parameter.
\textbf{Observation and design.} $d_y=4$, the effluent quality variables
$(S_S,S_{NH},S_{NO},S_O)$ with heteroscedastic noise
$\sigma=\sqrt{(0.05\,y)^2+0.5^2}$; designs are the aeration rate $k_La\in[10,300]$ and
the waste-sludge fraction $Q_w\in[0.01,0.2]$. The stiffness and the resulting adjoint
cost are what make policy training infeasible here.

\section{Multi-region SEIR}
\textbf{Forward model.} $R$ coupled regions, each with $(S_r,E_r,I_r,R_r)$ and
population $N=10^4$. With total infection $\textstyle\sum_r I_r$ and cross-region
pressure $\mathrm{cross}_r=\kappa(\sum_{r'}I_{r'}-I_r)$, and per-region designs
$(v_r,d_r,\tau_r)$ (vaccination, distancing, treatment),
\begin{align*}
\dot S_r &= -\beta_r(1-\varepsilon_d d_r)S_r(I_r+\mathrm{cross}_r)/N-v_r\varepsilon_v S_r, &
\dot E_r &= \beta_r(1-\varepsilon_d d_r)S_r(I_r+\mathrm{cross}_r)/N-\sigma E_r,\\
\dot I_r &= \sigma E_r-\gamma I_r-\tau_r\varepsilon_t I_r, &
\dot R_r &= \gamma I_r+\tau_r\varepsilon_t I_r+v_r\varepsilon_v S_r,
\end{align*}
from $S_r(0)=N-I_0$, $I_r(0)=I_0=10$, $E_r(0)=R_r(0)=0$.
\textbf{Prior.} $d_\theta=6+R$: shared $\sigma\sim\mathrm{LogNormal}(\log 0.2,0.3^2)$,
$\gamma\sim\mathrm{LogNormal}(\log 0.1,0.3^2)$, $\kappa\sim\mathrm{LogNormal}(\log 0.05,0.5^2)$;
per-region $\beta_r\sim\mathrm{LogNormal}(\log 0.4,0.5^2)$; and efficacies
$\varepsilon_v,\varepsilon_d,\varepsilon_t\sim\mathrm{Beta}(2,2)$.
\textbf{Observation and design.} $d_y=2R$ (exposed and infected per region) with
lognormal noise of scale $0.15$; $d_\xi=3R$ with $v_r\in[0,0.05]$, $d_r\in[0,1]$,
$\tau_r\in[0,0.5]$. We use $R=3$ as the main system ($d_\theta=9$); $R=5$ and $R=8$
appear in the saturation survey.

\begin{table}[h]
\centering\footnotesize
\caption[Experiment configurations]{Per-system experiment configurations for the
deployed PACE-ftk runs. Optimiser is forward-only cross-entropy (CEM); $L/J$ is the
contrastive budget ($J$ particles for PACE, $L_{\mathrm{train}}$ for DAD);
$N_{\mathrm{cand}}$ is the candidate count per CEM round; $C$ is the number of CEM
refinement rounds ($C{=}0$ is an unrefined random pool). The PACE-grad ablation uses the
same proposal and candidate budget with gradient refinement instead of CEM.}
\label{tab:config}
\begin{tabular}{lrrrlrrrr}
\toprule
System & $d_\theta$ & $d_\xi$ & $T$ & opt. & $L/J$ & $N_{\mathrm{cand}}$ & $C$ & $L_{\mathrm{eval}}$ \\
\midrule
CES & 4 & 6 & 10 & CEM & 1000 & 1000 & 1 & $10^7$ \\
ASM1 & 12 & 2 & 5 & CEM & 500 & 50 & 0 & $5{\times}10^4$ \\
SEIR R=3 & 9 & 9 & 10 & CEM & 1000 & 1000 & 1 & $10^7$ \\
Monod & 2 & 1 & 14 & CEM & 1000 & 1000 & 1 & $5{\times}10^3$ \\
Goodwin & 5 & 1 & 10 & CEM & 1000 & 1000 & 1 & $5{\times}10^3$ \\
Src.\ loc.\ $p{=}2$ & 4 & 2 & 10 & CEM & 1000 & 1000 & 1 & $10^5$ \\
Src.\ loc.\ $p{=}5$ & 10 & 5 & 30 & CEM & 1000 & 1000 & 1 & $10^5$ \\
Src.\ loc.\ $p{=}8$ & 16 & 8 & 30 & CEM & 1000 & 1000 & 1 & $10^5$ \\
PK 1-comp $T{=}3$ & 4 & 1 & 3 & CEM & 1000 & 1000 & 1 & $2{\times}10^3$ \\
PK 1-comp $T{=}5$ & 4 & 1 & 5 & CEM & 1000 & 1000 & 1 & $2{\times}10^3$ \\
PK 8-comp & 20 & 1 & 5 & CEM & 1000 & 1000 & 1 & $2{\times}10^3$ \\
PI3K & 6 & 2 & 8 & CEM & 1000 & 1000 & 1 & $10^5$ \\
\bottomrule
\end{tabular}
\end{table}

\section{Baseline policy (DAD) configuration}
The DAD baseline follows \citet{foster2021deep}: a permutation-invariant history encoder
(a shared per-step MLP applied to each $(\xi_i,y_i)$ pair, sum-pooled across steps)
feeding a design-emitter MLP that outputs the next design. It is trained by Adam
stochastic gradient ascent on the trajectory sPCE bound (Eq.~\eqref{eq:traj-spce}) with
$L_{\mathrm{train}}$ prior contrastives, $L_{\mathrm{train}}=1000$ for all systems except
CES ($1024$), with per-system gradient-step budgets of $3{,}000$ to $50{,}000$ and batch
sizes of $8$ to $128$; the resulting training times are in Table~\ref{tab:walltime}. For
SEIR R=3, $L_{\mathrm{train}}$ is additionally swept over
$\{10^2,10^3,10^4,5{\times}10^4,10^5\}$ (Table~\ref{tab:dad-sweep}). DAD is evaluated on
the same $\theta_{\mathrm{true}}$ seeds and at the same $L_{\mathrm{eval}}$ as PACE. On
CES, training diverges to NaN before convergence (the censored-likelihood instability
also reported by \citet{hedman2025stepdad}), and the reported result uses the best
checkpoint before divergence; on ASM1, and on SEIR R=3 at the budget needed to escape
saturation, training is infeasible (stiff-ODE adjoint with
$L_{\mathrm{escape},50}$ beyond any feasible $L_{\mathrm{train}}$;
Section~\ref{sec:m-forward}).

\chapter{Neural posterior estimator}
\label{app:npe}

\section{Data generation and posterior-target correctness}
Training data consist of $N$ triples $(\theta_i,\xi^{(i)}_{1:T},y^{(i)}_{1:T})$:
$\theta_i\sim p(\theta)$, designs $\xi^{(i)}_t\sim\mathrm{Uniform}(\Xi)$ sampled
independently for all $T$ steps, and observations
$y^{(i)}_t\sim p(y_t\mid\theta_i,\xi^{(i)}_t)$. Because the design distribution
$p(\xi_{1:t})$ does not depend on $\theta$, the training-joint posterior equals the
true posterior:
\begin{equation}
\label{eq:npe-correct}
p_{\mathrm{train}}(\theta\mid h_t)
= \frac{p(\theta)\,p(\xi_{1:t})\,p(y_{1:t}\mid\theta,\xi_{1:t})}{p_{\mathrm{train}}(h_t)}
= \frac{p(\theta)\,p(y_{1:t}\mid\theta,\xi_{1:t})}{p(y_{1:t}\mid\xi_{1:t})}
= p(\theta\mid y_{1:t},\xi_{1:t}),
\end{equation}
where $p(\xi_{1:t})$ cancels. The NPE therefore targets the correct posterior for
\emph{any} history $h_t$, regardless of how the designs were chosen. At deployment PACE
queries the NPE with histories containing optimally selected designs; the posterior is
determined by Bayes' rule given the data, not by the training design distribution. Each
training sample is a full $T$-step trajectory, and random truncation
($t_i\sim\mathrm{Uniform}\{1,\dots,T\}$, masking steps $t>t_i$) teaches the network to
produce posteriors from any prefix length.

\section{Set-equivariant conditioning}
The posterior $p(\theta\mid h_t)$ is permutation-invariant with respect to the elements
of $h_t=\{(y_s,\xi_s,s)\}_{s=1}^{t}$, and the conditioning architecture reflects this.
Each element is encoded as a triple $(y_s,\xi_s,s/T)$: a $z$-scored observation
(log-transformed under lognormal noise), a $z$-scored design, and a normalised time
index. A shared trial network $f_{\mathrm{trial}}:\mathbb{R}^{d_y+d_\xi+1}\to\mathbb{R}^{d_e}$
(2 layers, SiLU, $d_e=64$) processes each triple independently,
$e_s=f_{\mathrm{trial}}(y_s,\xi_s,s/T)$. The embeddings are sum-pooled,
$\bar e=\sum_{s=1}^{t}e_s$ (sum, not mean, to preserve the magnitude signal so that $t$
observations produce a different embedding than one), a count feature $t/T$ is appended,
and a set network $f_{\mathrm{set}}:\mathbb{R}^{d_e+1}\to\mathbb{R}^{d_{\mathrm{emb}}}$
(2 layers, SiLU, $d_{\mathrm{emb}}=128$) produces the final conditioning vector.

\section{Density estimator and training}
$q_\phi(\theta\mid h_t)$ is a neural spline flow~\citep{durkan2019neural} with $12$
coupling transforms and hidden dimension $256$. Parameters are mapped to an
unconstrained space (log for $(0,\infty)$, logit for $(0,1)$). NSF is chosen over a
flow-matching estimator because the importance reweighting requires fast, exact
evaluation of $\log q_\phi$ ($\sim$1\,ms, versus $\sim$500\,ms for the ODE solve of a
flow-matching posterior). Training uses the \texttt{sbi} library~\citep{tejerocantero2020sbi}
with \texttt{z\_score\_x=`none'}, because the flattened condition includes binary mask
channels that $z$-scoring would corrupt. The hyperparameters are:
\begin{center}\small
\begin{tabular}{ll}
\toprule
Hyperparameter & Value \\
\midrule
Coupling transforms / hidden dim & 12 / 256 \\
Learning rate / batch size & $5\times10^{-4}$ / 256 \\
Max epochs / early-stopping patience & 500 / 30 epochs \\
Trial network & 2 layers, 64 hidden, SiLU \\
Set network & 2 layers, 128 hidden, SiLU \\
\bottomrule
\end{tabular}
\end{center}
